\NeedsTeXFormat{LaTeX2e}
\documentclass[a4paper,10pt,reqno]{amsart}
%\documentclass[12pt]{amsart}
\usepackage{amscd}
\usepackage{amsmath}
\usepackage{ulem}
\usepackage{amssymb}
\usepackage{amsthm}
\usepackage[usenames,dvipsnames]{color}
\usepackage{color,soul}
\usepackage{enumerate, mdwlist}
%\usepackage{dcpic,pictexwd}
% Draft-only label-name display is omitted in this reading version.
\usepackage{anysize}
\usepackage{tikz}
\usepackage{url}
\usepackage{esint}
\usepackage{mathrsfs}

\input xy
\xyoption{all}

\begin{document}

%theorem environments
\newtheorem{thm}{Theorem}
\newtheorem{prop}[thm]{Proposition}
\newtheorem{conj}[thm]{Conjecture}
\newtheorem{lem}[thm]{Lemma}
\newtheorem{cor}[thm]{Corollary}
\newtheorem{axiom}[thm]{Axiom}
%\newtheorem{alpaca}{Sublemma}[llama] %subbits of llamas?
\newtheorem{sheep}[thm]{Corollary}
\newtheorem{deff}[thm]{Definition}
%\newtheorem{lett}{Step}[section] %steps in chapter2 section 3
\newtheorem{fact}[thm]{Fact}
%\newtheorem{assume}{Notation/Assumption}[section]
\newtheorem{example}[thm]{Example}
\newtheorem{assumption}[thm]{Assumption}
\newtheorem{remark}[thm]{Remark}
\newtheorem{quest}[thm]{Question}
\newtheorem{zample}[thm]{Example}
\newtheorem{hyp}[thm]{Hypothesis}
\newtheorem{claim}[thm]{Claim}

\newcommand{\sthat}{\hspace{.1cm}| \hspace{.1cm}}
%shortcuts for curve names
%\newcommand{\p}{\mathbb{P}^1}
%\newcommand{\ga}{\mathbb{G}_a}
%\newcommand{\gm}{\mathbb{G}_m}
%\newcommand{\spec}{ \mathit{Spec} }
\newcommand{\id}{\operatorname{id} }
\newcommand{\acl}{\operatorname{acl}}
\newcommand{\dcl}{\operatorname{dcl}}
\newcommand{\irr}{\operatorname{irr}}
\newcommand{\aut}{\operatorname{Aut}}
%\newcommand{\deg}{\operatorname{deg}}
\newcommand{\fix}{\operatorname{Fix}}
\newcommand{\abs}{\operatorname{abs}}
\newcommand{\ord}{\operatorname{ord}}
\newcommand{\GL}{\operatorname{GL}}
\newcommand{\PGL}{\operatorname{PGL}}
\newcommand{\ch}{\operatorname{ch}}

\newcommand{\oo}{\mathcal{O}}
\newcommand{\aaa}{\mathcal{A}}
\newcommand{\mm}{\mathcal{M}}
\newcommand{\curg}{\mathcal{G}}
\newcommand{\bbf}{\mathbb{F}}
\newcommand{\A}{\mathbb{A}}
\newcommand{\R}{\mathbb{R}}
\newcommand{\PP}{\mathbb{P}}
\newcommand{\Q}{\mathbb{Q}}
\newcommand{\C}{\mathbb{C}}
\newcommand{\cc}{\mathcal{C}}
\newcommand{\dd}{\mathcal{D}}
\newcommand{\N}{\mathbb{N}}
\newcommand{\Z}{\mathbb{Z}}
\newcommand{\cF}{\mathcal F}
\newcommand{\cB}{\mathcal B}
\newcommand{\cU}{\mathcal U}
\newcommand{\cV}{\mathcal V}
\newcommand{\cG}{\mathcal G}
\newcommand{\cD}{\mathcal D}
\newcommand{\curly}{\mathcal{C}}
\newcommand{\durly}{\mathcal{D}}
\newcommand{\fff}{\mathcal{F}}
\newcommand{\calc}{\mathcal{C}}
\newcommand{\calS}{\mathcal{S}}
\newcommand{\GG}{\mathbb{G}}
%\newcommand{\PP}{\mathbb{P}}
\newcommand{\Gal}{\mathrm{Gal}}
\newcommand{\Aut}{\mathrm{Aut}}
\newcommand{\signature}{\mathrm{sign}}
%\newcommand{\GL}{\mathrm{GL}}
\newcommand{\uu}{\underline}
\newcommand{\di}{}

\definecolor{mypink3}{cmyk}{0, 0.7808, 0.4429, 0.1412}


\newcommand{\lean}[1]{%
  \par\smallskip\begingroup\normalfont\sffamily\color{black}%
  \noindent\textbf{Lean implementation.} #1\par\endgroup\smallskip}

\DeclareRobustCommand{\hlgreen}[1]{{\sethlcolor{green}\hl{#1}}}
\DeclareRobustCommand{\hlcyan}[1]{{\sethlcolor{cyan}\hl{#1}}}
\newcommand{\Fmodtor}{F^\times / \mu(F)}
\newcommand{\ind}{\mathrm{index}\,}
\title[Roth's theorem]{Formalizing Roth's theorem}

\author{Ramin Takloo-Bighash}
\author{Samuel Yin}
\begin{abstract}
Here we put down the groundwork for formalizing Roth's theorem.
\end{abstract}
\maketitle

\section{Derivatives}

We start by defining differential operators in $p\geq1$ variables.
For a multi-index $i=(i_1,\dots,i_p)\in\N^p$, define the \emph{Hasse differential operator}
$\Delta_i$ by its effect on monomials:
\[
\Delta_i(x_1^{k_1}\cdots x_p^{k_p})=
\begin{cases}
\displaystyle\left(\prod_{j=1}^p {k_j\choose i_j}\right)
 x_1^{k_1-i_1}\cdots x_p^{k_p-i_p},& i_j\leq k_j\text{ for every }j,\\[3pt]
0,&\text{otherwise}.
\end{cases}
\]
We extend this operator linearly to all polynomials. It has order
$|i|=i_1+\cdots+i_p$. The definition works over $\Z$, $\Q$, or $\R$
(and over any commutative semiring by casting the binomial coefficients).
Over a $\Q$-algebra it agrees with
\[
\Delta_i=\frac{1}{i_1!\cdots i_p!}
 \left(\frac{\partial}{\partial x_1}\right)^{i_1}\cdots
 \left(\frac{\partial}{\partial x_p}\right)^{i_p}.
\]
When $i_j\leq k_j$, one may equivalently write
${k_j\choose i_j}={k_j\choose k_j-i_j}$.
Unless explicitly stated otherwise, a differential operator below means one
of these individual normalised operators, not an arbitrary linear combination.
This algebraic definition avoids the technicalities of Fr\'echet derivatives.

\lean{Represent polynomials by \texttt{MvPolynomial} and multi-indices
by finitely supported maps $\sigma\to_0\N$. Define each $\Delta_d$ as a
linear map and prove the coefficient identity
$[X^e]\Delta_dP=(\prod_j\binom{e_j+d_j}{d_j})[X^{e+d}]P$,
where $[X^e]$ denotes coefficient extraction; this shows that distinct
surviving monomials do not merge under differentiation. Use the univariate
\texttt{Polynomial.hasseDeriv}\footnote{\url{https://leanprover-community.github.io/mathlib4_docs/Mathlib/Algebra/Polynomial/HasseDeriv.html}}
as a model, and prove compatibility with coefficient maps before passing
between $\Z$, $\Q$ and $\R$.}

\section{Wronskians}

If $l\geq1$ and $\phi_0(t), \dots, \phi_{l-1}(t)$ are $l$ polynomials in the variable $t$, we call the determinant
$$
W(\phi_0, \dots, \phi_{l-1})(t) = \det \left(\frac{1}{\mu !} \frac{d^\mu}{d t^\mu} \phi_\nu(t)\right)_{0 \leq \mu \leq l-1 \atop 0 \leq \nu \leq l-1}
$$
the Wronskian of the polynomials.

\begin{lem}
    If $\phi_0(t), \dots, \phi_{l-1}(t)$ have rational coefficients, then the Wronskian vanishes if and only if they are linearly dependent, i.e., there are rational constants $c_0, \dots, c_{l-1}$, not all of which are zero, such that
    $$
    c_0 \phi_0(t) + \dots + c_{l-1} \phi_{l-1}(t) = 0.
    $$
\end{lem}
\begin{proof}
    Linear dependence gives a nontrivial constant linear relation between
    the columns, so the determinant vanishes.
    Conversely, suppose the polynomials are linearly independent. By invertible
    column operations over $\Q$, replace them by a basis $g_0,\dots,g_{l-1}$ of
    their span having distinct degrees $d_0<\cdots<d_{l-1}$.
    Such a basis is obtained by eliminating equal leading degrees: the
    successive quotients of the degree filtration have dimension at most one.
    If $A\in\mathrm{GL}_l(\Q)$ is this change-of-basis matrix, then
    $W(g_0,\dots,g_{l-1})=W(\phi_0,\dots,\phi_{l-1})\det A$.

    Write $a_\nu\ne0$ for the leading coefficient of $g_\nu$ and set
    $D=\sum_\nu d_\nu-l(l-1)/2\geq0$. The coefficient of $t^D$ in the new
    Wronskian is
    \[
    \left(\prod_{\nu=0}^{l-1}a_\nu\right)
    \det\left({d_\nu\choose\mu}\right)_{\mu,\nu=0}^{l-1}
    =\left(\prod_{\nu=0}^{l-1}a_\nu\right)
      \frac{\prod_{0\leq i<j<l}(d_j-d_i)}{\prod_{\mu=0}^{l-1}\mu!}\ne0.
    \]
    Indeed, $\binom{X}{\mu}$ has degree $\mu$ and leading coefficient
    $1/\mu!$, so its evaluation determinant is the Vandermonde determinant
    multiplied by $\prod_\mu1/\mu!$. All other contributions have lower total
    degree. Thus the new Wronskian, and hence the original one, is nonzero.
\end{proof}

\lean{Index both rows and columns by \texttt{Fin l}, the type of integers
$0,\ldots,l-1$, and take the determinant of the matrix with entries
$\Delta_\mu\phi_\nu$ in $\Q[t]$. Using the
matrix determinant API\footnote{\url{https://leanprover-community.github.io/mathlib4_docs/Mathlib/LinearAlgebra/Matrix/Determinant/Basic.html}},
split the independence criterion into change of basis, construction of a
basis with distinct degrees, and the Vandermonde leading-coefficient
calculation. This keeps the construction valid for arbitrary finite
families, rather than just two polynomials.}

Suppose $l\geq1$ and $\phi_0(x_1, \dots, x_p), \phi_1(x_1, \dots, x_p), \dots, \phi_{l-1}(x_1, \dots, x_p)$ are $l$ polynomials in the $p$ variables $x_1, \dots, x_p$. Also suppose $\Delta_0, \Delta_1, \dots, \Delta_{l-1}$ are $l$ differential operators such that for each $i$, $0 \leq i \leq l-1$, the order of $\Delta_i$ is at most $i$. We call the determinant
$$
W(\phi_0, \dots, \phi_{l-1}; \Delta_0, \dots, \Delta_{l-1}) = \det \left(\Delta_\mu \phi_\nu (x_1, \dots, x_p) \right)_{ 0 \leq \mu \leq l-1 \atop 0 \leq \nu \leq l-1}
$$
a generalized Wronskian. This is understood as a polynomial in the variables $x_1, \dots, x_p$.  If $p > 1$ and $l >1$ we have more than one generalized Wronskian.

\begin{lem}
    If $\phi_0, \dots, \phi_{l-1}$ are polynomials over a characteristic-zero field and are linearly dependent over that field, then all of their generalized Wronskians are identically zero.
\end{lem}
\begin{proof}
    Let $\sum_\nu c_\nu\phi_\nu=0$ be a nontrivial dependence relation
    over the coefficient field. Applying each $\Delta_\mu$ gives
    $\sum_\nu c_\nu\Delta_\mu\phi_\nu=0$. Choose $\nu_0$ with
    $c_{\nu_0}\ne0$ and express that column as a linear combination of the
    other columns. Multilinearity and alternation of the determinant then
    give zero.
\end{proof}

\begin{lem}\label{roth:lem1}
    If $\phi_0(x_1, \dots, x_p), \dots, \phi_{l-1}(x_1, \dots, x_p)$ are $l$ linearly independent polynomials in $p$ variables with rational coefficients, then at least one of their generalized Wronskians is not identically zero.
\end{lem}
\begin{proof}
    Let $k\geq2$ be an integer which is bigger than the degrees of the polynomials
    $\phi_0, \phi_1, \dots, \phi_{l-1}$ in each of the variables $x_1, \dots, x_p$. We will fix $k$ throughout this argument so we will often suppress the dependence of our objects on $k$. For each $\nu$, $0 \leq \nu \leq l-1$, let
    $$
    h_\nu(t) = \phi_\nu(t, t^k, t^{k^2}, \dots, t^{k^{p-1}})
    $$
    as polynomials of the variable $t$. We claim that these polynomials are linearly independent. For each $\nu$ write
    $$
    \phi_\nu(x_1, \dots, x_p) = \sum_{s_1=0}^{k-1} \dots \sum_{s_p = 0}^{k-1} b^{(\nu)}(s_1, \dots, s_p) x_1^{s_1} \cdots x_p^{s_p}.
    $$
    Suppose we have a linear dependence relation of the form
    $$
    \sum_{\nu=0}^{l-1} c_\nu h_\nu(t) = 0.
    $$
    We have
    $$
    \sum_{\nu=0}^{l-1} c_\nu \sum_{s_1=0}^{k-1} \dots \sum_{s_p = 0}^{k-1} b^{(\nu)}(s_1, \dots, s_p) t^{s_1 + k s_2 + \dots + k^{p-1} s_p} = 0.
    $$

FACT: Since each $s_i$ is strictly smaller than $k$, the representation of any natural number in the form
$$
s_1 + k s_2 + \dots + k^{p-1} s_p
$$
is unique.

Comparing the coefficient of $t^{s_1+ks_2+\cdots+k^{p-1}s_p}$ gives
\[
\sum_{\nu=0}^{l-1}c_\nu b^{(\nu)}(s_1,\dots,s_p)=0
\qquad(0\leq s_j<k).
\]
It follows, by collecting coefficients of the multivariable monomials, that
\[
\sum_{\nu=0}^{l-1}c_\nu\phi_\nu
 =\sum_{0\leq s_1,\dots,s_p<k}
   \left(\sum_{\nu=0}^{l-1}c_\nu b^{(\nu)}(s_1,\dots,s_p)\right)
   x_1^{s_1}\cdots x_p^{s_p}=0.
\]
The linear independence of the $\phi_\nu$ implies that every $c_\nu$ is zero.

\lean{Define Kronecker substitution as the algebra homomorphism
$x_j\mapsto t^{k^{j-1}}$, and prove injectivity only on the subspace where
every individual exponent is less than $k$. This restriction is essential:
when $p\geq2$, the nonzero polynomial $x_2-x_1^k$ maps to zero. Uniqueness
of base-$k$ expansions gives the restricted coefficient comparison, which
then transports the required linear independence.}

We conclude that the Wronskian
$$
W(h_0, \dots, h_{l-1})
$$
is not identically zero. We now relate this Wronskian to the Wronskians of the multivariable polynomials $\phi_\nu$.

Set
$$
\gamma(t) = (t, t^k, t^{k^2}, \dots, t^{k^{p-1}}).
$$
For a polynomial $f$ in $x_1, \dots, x_p$, the composite $f(\gamma(t))$ is a polynomial in the single variable $t$. We now relate its derivatives to the partial derivatives of $f$.

FACT: For each $\nu$,
$$
\frac{d}{dt}h_\nu(t) = \frac{\partial \phi_\nu} {\partial x_1 } (\gamma(t)) + k t^{k-1} \frac{\partial \phi_\nu} {\partial x_2 } (\gamma(t)) + \dots + k^{p-1}t^{k^{p-1}-1}\frac{\partial \phi_\nu} {\partial x_p } (\gamma(t)).
$$

This is the algebraic chain rule for polynomial substitution. The higher derivatives are handled by the uniform expansion below.

Then for each $\mu$, there is a natural number $r(\mu)$ depending only on $\mu$, polynomials $f^\mu_1(t), \dots, f^\mu_{r(\mu)}(t)$ with rational coefficients, and differential operators $\Delta_\mu^{(1)}, \dots, \Delta_\mu^{({r(\mu)})}$ of order less than or equal to $\mu$ such that for all $\nu$
$$
\frac{d^\mu}{dt^\mu} h_\nu(t) = f^\mu_1(t) (\Delta_\mu^{(1)}\phi_\nu)(\gamma(t)) + \dots + f^\mu_{r(\mu)}(t) (\Delta_\mu^{({r(\mu)})}\phi_\nu)(\gamma(t)).
$$

For completeness, this follows by induction on $\mu$, uniformly for all
input polynomials. For $\mu=0$, use the identity operator with coefficient $1$.
If a term at stage $\mu$ is $a(t)(\Delta_d\phi)(\gamma(t))$, its derivative is
\[
a'(t)(\Delta_d\phi)(\gamma(t))+
 a(t)\sum_{j=1}^p\gamma'_j(t)(d_j+1)
                 (\Delta_{d+e_j}\phi)(\gamma(t)).
\]
Here $\partial_j\Delta_d=(d_j+1)\Delta_{d+e_j}$ follows directly on
monomials. The orders are at most $\mu+1$, and all coefficients are rational
polynomials independent of $\phi$. Thus the same finite list works for every
$\phi_\nu$; $p$ and $k$ remain fixed throughout.
\lean{Treat $\phi\mapsto\phi(\gamma(t))$ as an algebra homomorphism and
prove the first-order chain rule using
algebraic partial derivatives\footnote{\url{https://leanprover-community.github.io/mathlib4_docs/Mathlib/Algebra/MvPolynomial/PDeriv.html}}
and polynomial induction, rather than analytic differentiation.
Represent the higher-order expansion by a finitely supported family of
coefficient polynomials indexed by $d$ with $|d|\leq\mu$; the recurrence
above makes these coefficients independent of $\phi$. Convert ordinary
iterated derivatives to Hasse derivatives by factorial scaling only over
$\Q$.}

Let
$$
\mathcal S = \prod_{\mu = 0}^{l-1} \{1, \dots, r(\mu)\}
$$
For each
$$\underline{i} = (i_0, \dots, i_{l-1}) \in
\mathcal S,
$$
set
$$
f_{\underline{i}}(t) = \prod_{\mu = 0}^{l-1} \frac{1}{\mu !} f^\mu_{i_\mu}(t).
$$
Then multilinearity of the determinant in each row gives
$$
W(h_0, \dots, h_{l-1}) = \det \left(\frac{1}{\mu !} \frac{d^\mu}{d t^\mu} h_\nu(t)\right)_{0 \leq \mu \leq l-1 \atop 0 \leq \nu \leq l-1}
$$
$$
=\det \left(\frac{1}{\mu !} f^\mu_1(t) (\Delta_\mu^{(1)}\phi_\nu)(\gamma(t)) + \dots + \frac{1}{\mu !}f^\mu_{r(\mu)}(t) (\Delta_\mu^{({r(\mu)})}\phi_\nu)(\gamma(t))\right)_{0 \leq \mu \leq l-1 \atop 0 \leq \nu \leq l-1}.
$$
$$
= \sum_{\underline i = (i_0, \dots, i_{l-1}) \in \mathcal S} f_{\underline{i}}(t) \det \left(\Delta_\mu^{(i_\mu)}\phi_\nu(\gamma(t)) \right)_{0 \leq \mu \leq l-1 \atop 0 \leq \nu \leq l-1}
$$
$$
= \sum_{\underline i = (i_0, \dots, i_{l-1}) \in \mathcal S} f_{\underline{i}}(t) W(\phi_0, \dots, \phi_{l-1}; \Delta_0^{(i_0)}, \Delta_1^{(i_1)}, \dots, \Delta_{l-1}^{(i_{l-1})})(\gamma(t)).
$$
Since $W(h_0, \dots, h_{l-1})$ is non-vanishing, this means there is a choice of differential operators such that $W(\phi_0, \dots, \phi_{l-1}; \Delta_0^{(i_0)}, \Delta_1^{(i_1)}, \dots, \Delta_{l-1}^{(i_{l-1})})$ is non-vanishing.
\end{proof}

\begin{lem}\label{roth:lem2}
    Let $R(x_1, \dots, x_p)$ be a polynomial in $p \geq 2$ variables with integral coefficients which is not identically zero. Let $R$ be of degree at most $r_j$ in $x_j$ for $j = 1, \dots, p$. Then there is an integer $l$ satisfying
    $$
    1 \leq l \leq r_p + 1
    $$
    and there are differential operators $\Delta_0, \dots, \Delta_{l-1}$ on the variables $x_1, \dots, x_{p-1}$ of orders at most $0, \dots, l-1$ respectively such that if
$$
F(x_1, \dots, x_p) = \det \left(
\Delta_\mu \frac{1}{\nu!}\left(\frac{\partial}{\partial x_p} \right)^\nu R
\right)_{0 \leq \mu \leq l-1 \atop 0 \leq \nu \leq l-1 }
$$
then
\begin{enumerate}
    \item $F$ has integral coefficients and it is not identically zero;
    \item we have
    $$
    F(x_1, \dots, x_p) = U(x_1, \dots, x_{p-1}) V(x_p),
    $$
    where $U, V$ are nonzero, have integral coefficients, and $U$ is of degree at most $lr_j$ in $x_j$ for $j=1, \dots, p-1$ and $V$ is of degree at most $lr_p$ in $x_p$.
\end{enumerate}
\end{lem}

\begin{proof}
    Consider all representations of $R$ in the form
    $$
    R(x_1, \dots, x_p) = \sum_{\nu=0}^{l-1} \phi_\nu(x_p) \psi_\nu(x_1, \dots, x_{p-1})
    $$
    with $\phi_\nu, \psi_\nu$ polynomials with rational coefficients such that $\phi_\nu$'s are of degree at most $r_p$ and the $\psi_\nu$'s of degree at most $r_j$ in $x_j$ for $1 \leq j \leq p-1$. It's possible to find such a representation. For example, we can let $l-1=r_p$ and $\phi_\nu(x_p) = x_p^\nu$ by collecting the terms in terms of the powers of $x_p$. This in particular shows that there is a choice of representation such that $1 \leq l \leq r_p +1$.

    Among all such representations, choose one with $l$ least.
If $l=1$, both factors are nonzero since $R\ne0$, so the independence
assertions below are immediate. For $l\geq2$, we claim
    $$
    \phi_0(x_p), \dots, \phi_{l-1}(x_p)
    $$
    are linearly independent. If not, then, after reindexing a nontrivial dependence relation, we can write
    $$
    \phi_{l-1} = d_0 \phi_0 + \dots + d_{l-2} \phi_{l-2}
    $$
    for rational numbers $d_0, \dots, d_{l-2}$. Then we have
    $$
    R = \phi_0(\psi_0 + d_0 \psi_{l-1}) + \dots + \phi_{l-2}(\psi_{l-2} + d_{l-2} \psi_{l-1})
    $$
    which contradicts the minimality of $l$. Similarly,
    $$
    \psi_0(x_1, \dots, x_{p-1}), \dots, \psi_{l-1}(x_1, \dots, x_{p-1})
    $$
    are also linearly independent.

Let $W(x_p)$ be the Wronskian of the single variable polynomials $\phi_0(x_p), \dots, \phi_{l-1}(x_p)$. Then $W(x_p)$ is a polynomial with rational coefficients which is not identically zero. Let $G(x_1, \dots, x_{p-1})$ be some generalized Wronskian of
$$
    \psi_0(x_1, \dots, x_{p-1}), \dots, \psi_{l-1}(x_1, \dots, x_{p-1})
    $$
    which is not identically zero. Then
    $$
    W(x_p) = \det \left( \frac{1}{\mu !} \frac{d^\mu}{dx_p^\mu}\phi_\nu(x_p) \right)_{0 \leq \mu \leq l-1 \atop 0 \leq \nu \leq l-1}
    $$
    and
    $$
    G(x_1, \dots, x_{p-1}) = \det \left(\Delta_\mu \psi_\nu (x_1, \dots, x_{p-1}) \right)_{0 \leq \mu \leq l-1 \atop 0 \leq \nu \leq l-1}
    $$
    where $\Delta_0, \dots, \Delta_{l-1}$ are certain differential operators in the variables $x_1, \dots, x_{p-1}$ of orders at most $0, 1, \dots, l-1$, respectively. Let $A_{\mu\rho}=\Delta_\mu\psi_\rho$ and
$B_{\nu\rho}=\Delta_\nu\phi_\rho$, viewing both matrices in the common
polynomial ring. The matrix defining $F$ is $AB^{\mathsf T}$, so
$F=\det A\det B$ by multiplicativity and invariance
under transpose. Explicitly, we obtain
    $$
    GW = \det \left(\sum_{\rho =0}^{l-1} \Delta_\mu \frac{1}{\nu !} \frac{\partial^\nu}{\partial x_p^\nu} \phi_\rho (x_p) \psi_\rho(x_1, \dots, x_{p-1})\right)_{0 \leq \mu \leq l-1 \atop 0 \leq \nu \leq l-1}
    $$
    $$
    = \det \left( \Delta_\mu \frac{1}{\nu !} \frac{\partial^\nu}{\partial x_p^\nu} R \right)_{0 \leq \mu \leq l-1 \atop 0 \leq \nu \leq l-1}.
    $$

    If we let $F = GW$, then $F$ is expressible in the form required by the lemma. Since $R$ has integral coefficients, so does $F$. Choose $g\in\Q^\times$ so that $U=gG$ is an integral primitive
polynomial, by clearing denominators and then dividing by the gcd of its
coefficients. Put $V=g^{-1}W$. For a coefficient $v_b$ of $V$ and any
coefficient $u_a$ of $U$, the product $u_av_b$ is a coefficient of $F$:
the variables of the two factors are disjoint. It is therefore an integer.
Primitivity gives a finite integral linear combination $\sum_a z_au_a=1$,
so $v_b=\sum_a z_a(u_av_b)$ is an integer. Thus $U,V$ are integral and
nonzero, as required.
\lean{Construct the separated decomposition over $\Q$ from a basis of the
span of the coefficient polynomials of $R$ in $x_p$. Isolate the integral
rescaling as a separate lemma: if integer coefficients $u_a$ have gcd $1$
and every product $u_av_b$ is integral, then each rational $v_b$ is integral
by a B\'ezout relation. Keep $U\ne0$ and $V\ne0$ explicit, since the degree
comparisons need both facts.}

    Finally, since the variables are disjoint and both factors are nonzero, the individual degrees of $U,V$ equal the corresponding degrees of $F$. Since $F$ is a determinant of an $l\times l$ matrix whose entries are polynomials of degree at most $r_j$ in $x_j$, $1 \leq j \leq p$, $F$ will be of degree at most $l r_j$ in each variable $x_j$.
\end{proof}

\begin{lem}\label{roth:lem3}
Let $R$ and $F$ be as above, and suppose all the coefficients of $R$ have absolute value bounded by $B$. Then all the coefficients   of $F$ have absolute value bounded by
$$
((r_1 + 1) \cdots (r_p+1))^l l! B^l 2^{(r_1 + \dots + r_p)l}.
$$
\end{lem}
\begin{proof}
For any polynomial $f=\sum_e a_eX^e$, put
\[
abc(f)=\max_e|a_e|\quad(abc(0)=0),\qquad
\|f\|_1=\sum_e|a_e|.
\]
We prove the stronger bound for $\|F\|_1$. The triangle inequality gives
$\|f+kg\|_1\leq\|f\|_1+|k|\|g\|_1$. For monomials (including the zero
polynomial), $\|f_1\cdots f_l\|_1=\prod_\nu\|f_\nu\|_1$.

Set $S=r_1+\cdots+r_p$ and $N=\prod_j(r_j+1)$. Pad the monomial expansion
of $R$ by zero terms and write $R=\sum_{i=1}^N\alpha_i$, each
$\alpha_i=u_iX^{s(i)}$ with $|u_i|\leq B$ and $s(i)_j\leq r_j$.
An individual Hasse derivative of such a term is zero or a monomial whose
coefficient has absolute value at most
\[
|u_i|\prod_j{s(i)_j\choose d_j}\leq B2^S.
\]
This applies to $\Delta_\mu\Delta^{(p)}_\nu\alpha_i$, where
$\Delta^{(p)}_\nu=(1/\nu!)\partial_{x_p}^{\nu}$, since $\Delta_\mu$ uses
only the first $p-1$ variables.

Let $\mathcal X$ be the set of maps $\{0,\dots,l-1\}\to\{1,\dots,N\}$.
Multilinearity in the columns gives
\[
F=\sum_{f\in\mathcal X}
 \det\bigl(\Delta_\mu\Delta^{(p)}_\nu\alpha_{f(\nu)}\bigr)_{\mu,\nu=0}^{l-1}.
\]
Write $S_l$ for the permutations of $\{1,\dots,l\}$.
For $\sigma\in S_l$, the map $\nu\mapsto\sigma(\nu+1)-1$ permutes
$\{0,\dots,l-1\}$, matching the row and column indices of the matrix.
The signs have absolute value $1$. Thus, for each $f$,
\begin{align*}
\left\|\det\bigl(\Delta_\mu\Delta^{(p)}_\nu\alpha_{f(\nu)}\bigr)\right\|_1
&\leq\sum_{\sigma\in S_l}
 \left\|\prod_{\nu=0}^{l-1}
      \Delta_{\sigma(\nu+1)-1}\Delta^{(p)}_\nu\alpha_{f(\nu)}\right\|_1\\
&=\sum_{\sigma\in S_l}\prod_{\nu=0}^{l-1}
 \left\|\Delta_{\sigma(\nu+1)-1}\Delta^{(p)}_\nu\alpha_{f(\nu)}\right\|_1\\
&\leq l!(B2^S)^l.
\end{align*}
There are $N^l$ maps in $\mathcal X$. Hence
\[
abc(F)\leq\|F\|_1\leq N^ll!(B2^S)^l,
\]
which proves both the assertion and the promised stronger estimate.
\end{proof}

\lean{Define the coefficient maximum $abc(f)$ and the coefficient sum
$\|f\|_1$ separately using finite support, with both equal to zero at the
zero polynomial. Prove $abc(f)\leq\|f\|_1$, the triangle inequality and
submultiplicativity for $\|\cdot\|_1$, then use these for the determinant
bound. Keep the sharper derivative estimate separate: each output
coefficient of a fixed Hasse derivative receives at most one input
coefficient.}

\section{Index}
We start with the Taylor expansion of a polynomial in several variables.
\begin{thm}[Taylor's expansion for polynomials in several variables]
    Let $P(x_1, \dots, x_p)$ be a polynomial with real coefficients in the variables $x_1, \dots, x_p$ of degree at most $r_i$ in $x_i$ for $1 \leq i \leq p$. For each $1 \leq i \leq p$, pick a natural number $d_i \geq r_i$. Let $(\alpha_1, \dots, \alpha_p) \in \R^p$. Then
    $$
    P(x_1, \dots, x_p) = \sum_{i_1=0}^{d_1} \cdots \sum_{i_p=0}^{d_p} (\Delta_{i_1, \dots, i_p} P)(\alpha_1, \dots, \alpha_p) \prod_{j=1}^p ( x_j - \alpha_j)^{i_j}.
    $$
\end{thm}
\begin{proof}
    For a polynomial $P$ let
    $$
    T_{d_1, \dots, d_p} P(x_1, \dots, x_p) = \sum_{i_1=0}^{d_1} \cdots \sum_{i_p=0}^{d_p} (\Delta_{i_1, \dots, i_p} P)(\alpha_1, \dots, \alpha_p) \prod_{j=1}^p ( x_j - \alpha_j)^{i_j}
    $$
It is clear that
$$
T_{d_1, \dots, d_p} (P_1 + P_2) = T_{d_1, \dots, d_p} (P_1) + T_{d_1, \dots, d_p}(P_2)
$$
and for any constant $c$,
$$
T_{d_1, \dots, d_p}(c P) = c T_{d_1, \dots, d_p}(P).
$$
So it suffices to prove the theorem for monomials. For any tuple
$0\leq s_i\leq r_i\leq d_i$, distribution gives
$$
T_{d_1,\dots,d_p}(x_1^{s_1}\cdots x_p^{s_p})
 =\prod_{i=1}^pT_{d_i}(x_i^{s_i})=\prod_{i=1}^px_i^{s_i}
$$
by the one-variable binomial expansion $x_i^{s_i}=((x_i-\alpha_i)+\alpha_i)^{s_i}$.
\end{proof}

Let $P(x_1, \dots, x_p)$ be a nonzero polynomial with real coefficients in $p$ variables. Let $\alpha_1, \dots, \alpha_p$ be any real numbers. Let $r_1, \dots, r_p$ be any positive numbers. We define the {\it index} $\theta$ of $P$ at the point $(\alpha_1, \dots, \alpha_p)$ relative to $r_1, \dots, r_p$ as follows: We define $\theta$ to be the minimum of all numbers of the form
$$
\frac{i_1}{r_1} + \cdots + \frac{i_p}{r_p}
$$
where the minimum is over all tuples $(i_1, \dots, i_p)$ such that
$$
\Delta_{i_1, \dots, i_p} P(\alpha_1, \dots, \alpha_p) \ne 0.
$$
When we want to keep track of all data we write
$$
\ind(P; \alpha_1, \dots, \alpha_p; r_1, \dots, r_p).
$$
When there is no confusion about the data, we simply write $\ind(P)$.
Taylor's theorem shows that at least one derivative at the chosen point is
nonzero; only finitely many derivative multi-indices can be nonzero. Thus
this is an attained minimum. We additionally set $\ind(0)=+\infty$.
All comparisons involving zero polynomials below use this convention.
\lean{Let $P_\alpha(X)=P(X+\alpha)$, whose coefficient at $d$ is
$(\Delta_dP)(\alpha)$. Define the index as the minimum of
$\sum_jd_j/r_j$ over its nonzero coefficients, with values in
$[0,+\infty]$ and value $+\infty$ for the zero polynomial. For positive
integral $r_j$, an alternative is to set $D=\prod_jr_j$ and use the
natural weights $w_j=D/r_j$ in
\texttt{MvPowerSeries.weightedOrder}\footnote{\url{https://leanprover-community.github.io/mathlib4_docs/Mathlib/RingTheory/MvPowerSeries/Order.html}},
then divide the finite value by $D$ after embedding $P_\alpha$ as a power
series. Arbitrary positive real $r_j$ require the direct finite-support
minimum; do not convert an infinite order to a natural number.}

\begin{lem}
   We have $\ind (P) \geq 0$ and $\ind(P) = 0$ if and only if
$$
P(\alpha_1, \dots, \alpha_p) \ne 0.
$$
\end{lem}
\begin{proof}
    Every weight $\sum_j i_j/r_j$ is nonnegative. Since each $r_j>0$,
    such a weight is zero exactly when every $i_j=0$. The minimum is attained,
    and $\Delta_0P=P$, proving the equivalence. For $P=0$ the index is
    infinite, so the assertion remains valid.
\end{proof}

\begin{lem}\label{lem:index-of-derivative}
    Suppose $k_1, \dots, k_p$ are non-negative integers such that $\Delta_{k_1, \dots, k_p} P$ is not identically zero. Then
    $$
    \ind(\Delta_{k_1, \dots, k_p} P; \alpha_1, \dots, \alpha_p; r_1, \dots, r_p) \geq \ind( P; \alpha_1, \dots, \alpha_p; r_1, \dots, r_p) - \left(\frac{k_1}{r_1}  + \cdots + \frac{k_p}{r_p}\right).
    $$
\end{lem}
\begin{proof}
    Choose a multi-index $i$ attaining the index of $\Delta_kP$.
    The Hasse composition identity gives
    \[
    \Delta_i\Delta_kP=
       \left(\prod_{j=1}^p{i_j+k_j\choose i_j}\right)\Delta_{i+k}P.
    \]
    Since its value at $\alpha$ is nonzero, $(\Delta_{i+k}P)(\alpha)\ne0$.
    Therefore
    $\ind(P)\leq\sum_j(i_j+k_j)/r_j
       =\ind(\Delta_kP)+\sum_jk_j/r_j$, as required.
\end{proof}
\begin{lem}\label{roth:lem4-1}
    Let $P,Q$ be polynomials with real coefficients. With $\ind(0)=+\infty$, we have
    $$
    \ind(P+Q) \geq \min(\ind P, \ind Q),
    $$
    $$
    \ind(PQ) = \ind(P) + \ind(Q)
    $$
    if the indices are calculated relative to the tuples $(\alpha_1, \dots, \alpha_p)$ and $(r_1, \dots, r_p)$ as above.
\end{lem}
\begin{proof}
    Write $P(X+\alpha)=\sum_i a_iX^i$ and
    $Q(X+\alpha)=\sum_i b_iX^i$, and give $X^i$ weight
    $w(i)=\sum_j i_j/r_j$. All coefficients of $P+Q$ below the smaller
    index vanish, proving the addition inequality, including the case
    $P+Q=0$.

    For multiplication, if either factor is zero, both sides are $+\infty$. Otherwise let $u=\ind(P)$ and $v=\ind(Q)$, and let $P_u,Q_v$
    be the sums of terms of weights exactly $u,v$ in the translated
    polynomials. They are nonzero. Since the polynomial ring over $\R$ is
    an integral domain, $P_uQ_v\ne0$. All its terms have weight $u+v$,
    while every term involving a higher-weight component has weight
    strictly greater than $u+v$. Thus the lowest nonzero weight of the
    product is exactly $u+v$.
\end{proof}
\begin{lem}\label{roth:lem4-2}
    Let $P$ be a polynomial in the variables $x_1, \dots, x_{p-1}$ and $Q$ a polynomial in the variable $x_p$, and suppose $P, Q$ are not identically zero. Then for tuples $(\alpha_1, \dots, \alpha_p)$ and $(r_1, \dots, r_p)$ as above we have
    $$
\ind(PQ; \alpha_1, \dots, \alpha_p; r_1, \dots, r_p) = \ind(P; \alpha_1, \dots, \alpha_{p-1}; r_1, \dots, r_{p-1}) + \ind(Q; \alpha_p; r_p).
$$
\end{lem}
\begin{proof}
    When $P$ is embedded into the larger polynomial ring, its translated
    monomials have zero exponent in $x_p$, so its index is unchanged.
    Similarly, the index of the embedded $Q$ is its one-variable index.
    Apply the product formula of Lemma \ref{roth:lem4-1}.
\end{proof}

\lean{For integral $r_j$, apply
\texttt{weightedOrder\_mul}\footnote{\url{https://leanprover-community.github.io/mathlib4_docs/Mathlib/RingTheory/MvPowerSeries/NoZeroDivisors.html}}
through translation and the normalisation just described, under a
no-zero-divisors assumption. For arbitrary positive real $r_j$, formalise
the lowest-weight-part argument above instead. Prove invariance under
nonzero scalars and variable renaming, the disjoint-variable formula, and
scaling of the index weights as separate lemmas, so later proofs do not
repeat the conversions.}

For nonzero $c$ and $l>0$, the same definition gives
\[
\ind(cP;\alpha;r)=\ind(P;\alpha;r),\qquad
\ind(P;\alpha;lr)=\frac1l\ind(P;\alpha;r),
\]
with the second identity interpreted in the extended nonnegative reals.

\section{Index across sets of polynomials}

Let $m$ be a natural number $m \geq 1$. For positive integers $r_1, \dots, r_m$ and a real number $B \geq 1$ we let
$$
\mathcal R_m = \mathcal R_m (B; r_1, \dots, r_m)
$$
be the set of polynomials $R(x_1, \dots, x_m)$ in the variables $x_1, \dots, x_m$ such that:
\begin{enumerate}
    \item $R$ has integral coefficients and is not identically zero;
    \item $R$ has degree at most $r_j$ in $x_j$ for $1 \leq j \leq m$;
    \item the coefficients of $R$ have absolute values not exceeding $B$.
\end{enumerate}

Let $q_1, \dots, q_m$ denote positive integers. We define
$$
\Theta_m(B; q_1, \dots, q_m; r_1, \dots, r_m)
$$
to be the supremum of the numbers
$$
\ind (R; \frac{h_1}{q_1}, \dots, \frac{h_m}{q_m}; r_1, \dots, r_m)
$$
over all $R \in \mathcal R_m(B; r_1, \dots, r_m)$ and integers $h_i$, $1 \leq i \leq m$, with $\gcd(h_i, q_i)=1$.
This supremum is well-defined and satisfies $0\leq\Theta_m\leq m$:
the constant polynomial $1$ with every $h_i=1$ is allowed, and Taylor's theorem gives a nonzero
derivative with $0\leq i_j\leq r_j$, of weight at most $m$.
Also, $\Theta_m$ is nondecreasing in $B$, since increasing $B$ enlarges the
set of allowed polynomials.

\begin{lem}\label{roth:lem5}
    Suppose additionally that $q_1>1$. Then we have
    $$
    \Theta_1(B; q_1; r_1) \leq \frac{\log B}{r_1 \log q_1}.
    $$
\end{lem}
\begin{proof}
    Let $R\in\mathcal R_1(B;r_1)$ and $\gcd(h_1,q_1)=1$.
    Let $t\in\N$ be the multiplicity of $h_1/q_1$ as a root of $R$;
    Taylor's theorem gives $\theta=\ind(R;h_1/q_1;r_1)=t/r_1$.
    The polynomial $(q_1x_1-h_1)^t$ is primitive and divides $R$ over $\Q$.
    Gauss's lemma therefore gives a nonzero $Q\in\Z[x_1]$ with
    \[
    R(x_1)=(q_1x_1-h_1)^tQ(x_1).
    \]
    Comparing nonzero leading coefficients yields
    \[
    B\geq|\operatorname{lc}(R)|
       =q_1^t|\operatorname{lc}(Q)|\geq q_1^t.
    \]
    Since $\log q_1>0$ and $r_1>0$, taking logarithms gives
    $\theta\leq\log B/(r_1\log q_1)$. Taking the supremum proves the claim.
\end{proof}

\lean{First identify the one-variable index with $t/r_1$, where $t$ is the
natural-number root multiplicity. Use
Gauss's lemma\footnote{\url{https://leanprover-community.github.io/mathlib4_docs/Mathlib/RingTheory/Polynomial/GaussLemma.html}}
for the primitive polynomial $(q_1X-h_1)^t$, compare leading coefficients
in $\Z$, and only then cast to $\R$ and take logarithms, using $q_1>1$.
For $\Theta_m$, record nonemptiness and boundedness of the defining family,
together with lemmas for bounding an individual index and monotonicity in
$B$, before using the supremum.}

\begin{lem}\label{roth:lem6}
    Let $p \geq 2$ be a positive integer, and let $0 < \delta <1$. Let $r_1, \dots, r_p$ be positive integers satisfying
    $$
    r_p > 10 \delta^{-1}, \quad r_{j-1}/r_j > \delta^{-1}, \text{ for } 2 \leq j \leq p.
    $$
    Also, let $q_1, \dots, q_p$ be positive integers. Then
    \begin{equation}\label{roth:eq13}
    \Theta_p(B; q_1, \dots, q_p; r_1, \dots, r_p) \leq
    2 \max_l (\Phi + \Phi^{1/2}+ \delta^{1/2}),
    \end{equation}
    where the maximum is taken over integers $l$ satisfying
    $$
    1 \leq l \leq r_p +1,
    $$
    and where
    \begin{equation}\label{roth:eq15}
    \Phi = \Theta_1(M; q_p; lr_p) + \Theta_{p-1}(M; q_1, \dots, q_{p-1}; lr_1, \dots, lr_{p-1})
    \end{equation}
    and
    \begin{equation}\label{roth:eq16}
    M= (r_1+1)^{pl} l! B^l 2^{plr_1}.
    \end{equation}
\end{lem}
\begin{proof}
    Let $R\in\mathcal R_p(B;r_1,\dots,r_p)$ and choose integers $h_1,\dots,h_p$ with $\gcd(h_i,q_i)=1$ for every $i$. We need to show that the index
    $$
    \ind(R; \frac{h_1}{q_1}, \dots, \frac{h_p}{q_p}; r_1, \dots, r_p)
    $$
    is bounded by the right hand side of \eqref{roth:eq13}.

    The polynomial $R$ satisfies the conditions of Lemma \ref{roth:lem2}, and therefore there is an integer $l$ satisfying $1 \leq l \leq r_p+1$ and a polynomial $F$ as in Lemma \ref{roth:lem2}. By Lemma \ref{roth:lem3} the coefficients of $F$ have absolute value not exceeding
    $$
    ((r_1+1)\cdots (r_p+1))^l l! B^l 2^{(r_1+\cdots+r_p)l} < (r_1+1)^{pl}l! B^l 2^{pr_1 l} = M,
    $$
    as $r_1 > r_2 > \dots > r_p$. Since
    $$
    F = U(x_1, \dots, x_{p-1})V(x_p)
    $$
    and $U,V$ are nonzero integral polynomials in disjoint variables,
each coefficient of $F$ is a unique product of one coefficient from each
factor. Thus $abc(F)=abc(U)abc(V)$, while $abc(U),abc(V)\geq1$, and hence
    $$
    abc(U) < M, \quad abc(V) < M.
    $$

    By Lemma \ref{roth:lem2},
    $$
    U \in \mathcal{R}_{p-1}(M; lr_1, \dots, lr_{p-1}).
    $$
    Hence,
    $$
    \ind(U; \frac{h_1}{q_1}, \dots, \frac{h_{p-1}}{q_{p-1}}; l r_1, \dots, l r_{p-1}) \leq \Theta_{p-1}(M; q_1, \dots, q_{p-1}; lr_1, \dots, lr_{p-1}).
    $$
    Consequently, by scaling the weights,
    \[
    \ind(U;h_1/q_1,\dots,h_{p-1}/q_{p-1};r_1,\dots,r_{p-1})
    \leq l\Theta_{p-1}(M;q_1,\dots,q_{p-1};lr_1,\dots,lr_{p-1}).
    \]
    Similarly, $V(x_p)\in\mathcal R_1(M;lr_p)$, and
    $$
    \ind (V; \frac{h_p}{q_p}; r_p) \leq l \Theta_1(M; q_p; lr_p).
    $$
    By Lemma \ref{roth:lem4-2}, writing $a=(h_1/q_1,\dots,h_p/q_p)$,
    \begin{align*}
    \ind(F;a;r)
    &=\ind(U;a_1,\dots,a_{p-1};r_1,\dots,r_{p-1})
        +\ind(V;a_p;r_p)\\
    &\leq l\Theta_{p-1}(M;q_1,\dots,q_{p-1};lr_1,\dots,lr_{p-1})
       +l\Theta_1(M;q_p;lr_p)=l\Phi.
    \end{align*}
    Thus we have established
    \begin{equation}\label{roth:eq17}
        \ind F\leq l\Phi.
    \end{equation}

\lean{Keep a polynomial's degree bounds separate from the positive
parameters used to measure its index. Multiplying the index parameters
by $l>0$ divides the index by $l$, but does not change which polynomials
satisfy a degree bound. Here $U$ has degree bounds $lr_j$, so apply the
index-scaling identity to $U$ first and retain those bounds in
$\Theta_{p-1}(M;q_1,\ldots,q_{p-1};lr_1,\ldots,lr_{p-1})$.}

    Next, we find a lower bound for the index of $F$ in terms of the index of $R$.

    We start with an observation.
    Let
$$
\theta = \ind(R; \frac{h_1}{q_1}, \dots, \frac{h_p}{q_p}; r_1, \dots, r_p).
$$
    Let $i_1, \dots, i_{p-1}$ be non-negative integers such that $w=i_1+\cdots+i_{p-1}\leq l-1$.  If
    $$
    \Delta_{i_1, \dots, i_{p-1}} \frac{1}{\nu !} \frac{\partial^\nu}{\partial x_p^\nu} R(x_1, \dots, x_p)
    $$
    is not identically zero, then
    $$
    \ind (\Delta_{i_1, \dots, i_{p-1}} \frac{1}{\nu !} \frac{\partial^\nu}{\partial x_p^\nu} R(x_1, \dots, x_p); \frac{h_1}{q_1}, \dots, \frac{h_p}{q_p}; r_1, \dots, r_p)
    $$
    $$
    \geq \theta - \frac{i_1}{r_1}- \cdots - \frac{i_{p-1}}{r_{p-1}}- \frac{\nu}{r_p}
    $$
    $$
    \geq \theta - \frac{w}{r_{p-1}}- \frac{\nu}{r_p}.
    $$
    We have
    $$
    \frac{w}{r_{p-1}} \leq \frac{l-1}{r_{p-1}} \leq \frac{r_p}{r_{p-1}} < \delta.
    $$
    Since the index is never negative, the index is at least
    $$
    \max(0, \theta - \nu/r_p - \delta) \geq \max(0, \theta - \nu/r_p) - \delta.
    $$
    We have
    $$
    F = \det \left( \Delta_\mu \frac{1}{\nu !} \frac{\partial^\nu}{\partial x_p^\nu} R \right)_{0 \leq \mu \leq l-1 \atop 0 \leq \nu \leq l-1}.
    $$
    $$
    = \sum_{\sigma \in S_l} \signature(\sigma) \prod_{\nu=0}^{l-1} \Delta_{\sigma(\nu+1)-1} \frac{1}{\nu !} \frac{\partial^\nu}{\partial x_p^\nu} R.
    $$
    Here $S_l$ and the shifted indices have the same meaning as in the coefficient bound.
    We now estimate the index of $F$. Zero entries and zero permutation
    products have index $+\infty$ and cannot worsen the lower bound.
    Multiplication by a sign does not change the index. By Lemma \ref{roth:lem4-1},
    $$
    \ind(F) \geq \min_{\sigma \in S_l} \ind \prod_{\nu=0}^{l-1} \Delta_{\sigma(\nu+1)-1} \frac{1}{\nu !} \frac{\partial^\nu}{\partial x_p^\nu} R
    $$
    $$
    \geq \min_{\sigma \in S_l} \sum_{\nu=0}^{l-1} \ind \Delta_{\sigma(\nu+1)-1} \frac{1}{\nu !} \frac{\partial^\nu}{\partial x_p^\nu} R
    $$
    $$
    \geq \min_{\sigma \in S_l} \sum_{\nu=0}^{l-1} \max(0, \theta - \nu/r_p) - l\delta
    $$
    $$
    = \sum_{\nu=0}^{l-1} \max(0, \theta - \nu/r_p) - l\delta.
    $$
    We have proved
    $$
    \ind(F) \geq \sum_{\nu=0}^{l-1} \max(0, \theta - \nu/r_p) - l\delta.
    $$
    We can suppose $\theta r_p > 10$, for if not, then we have
    $$
    \theta \leq 10 r_p^{-1} < \delta < 2 \delta^{1/2},
    $$
    and this shows the inequality in the statement of the lemma is valid.

    Let $\theta r_p <l$. Then in order for
    $$
    \max(0, \theta - \nu/r_p)
    $$
    to be positive, we need $\nu<\theta r_p$. We may include the zero endpoint and sum over $\nu\leq\theta r_p$.
    \begin{equation}\label{roth:nonumber}
    \sum_{\nu=0}^{l-1} \max(0, \theta - \nu/r_p) = \sum_{0 \leq \nu \leq \theta r_p}(\theta  - \nu/r_p) = r_p^{-1}\sum_{0 \leq \nu \leq \theta r_p}(\theta r_p - \nu).
    \end{equation}

FACT: For any $\alpha \geq 0$ we have
$$
\sum_{0 \leq \nu \leq \alpha} (\alpha - \nu) \geq \frac{1}{2}[\alpha]^2.
$$
   Proof of fact:
$$
\sum_{0 \leq \nu \leq \alpha} (\alpha - \nu) = \alpha([\alpha]+1) - \sum_{0 \leq \nu \leq [\alpha]} \nu
$$
$$
= \alpha([\alpha]+1)- \frac{[\alpha]([\alpha]+1)}{2}
$$
$$
= ([\alpha]+1)(\alpha - \frac{[\alpha]}{2}) = ([\alpha]+1)([\alpha] + \{\alpha \} - \frac{[\alpha]}{2}) \quad (\text{here } \alpha = [\alpha] + \{\alpha \})
$$
$$
= ([\alpha]+1)( \frac{[\alpha]}{2}+ \{\alpha \} )
$$
$$
\geq [\alpha] \cdot \frac{[\alpha]}{2}
$$
as claimed.

FACT: If $\alpha >10$, then
$$
\frac{[\alpha]^2}{2} > \frac{1}{3}\alpha^2.
$$

Proof of fact: put $n=[\alpha]$. Since $\alpha>10$, we have $n\geq10$
and $\alpha<n+1\leq(11/10)n$. Therefore
\[
\alpha^2<\frac{121}{100}n^2<\frac32n^2,
\]
which is equivalent to $n^2/2>\alpha^2/3$.

   By applying these two facts we see that the expression  \eqref{roth:nonumber} is bounded below by
    $$
    \geq \frac{1}{2}r_p^{-1} [\theta r_p]^2
    $$
    $$
    > \frac{1}{3}r_p \theta^2.
    $$
    If $\theta r_p \geq l$,
    $$
    \sum_{\nu=0}^{l-1} \max(0, \theta - \nu/r_p) = \sum_{\nu=0}^{l-1} (\theta - \nu/r_p) \geq \frac{1}{2}l \theta.
    $$
In order to see this last inequality we proceed as follows:
$$
\sum_{\nu=0}^{l-1} (\theta - \nu/r_p) = l \theta - \frac{1}{r_p} \sum_{\nu = 0}^{l-1} \nu = l \theta - \frac{1}{r_p} \cdot \frac{l(l-1)}{2} \geq l \theta - \frac{1}{r_p} \cdot \frac{l^2}{2} \geq l\theta - \frac{1}{2}l \theta = \frac{1}{2}l \theta.
$$
Here we have used $- \frac{l}{r_p} \geq - \theta$.

    Combining everything, we obtain
    $$
    \ind F \geq \min (\frac{1}{2}l \theta, \frac{1}{3}r_p \theta^2) - l \delta.
    $$

\lean{Extract the bound for
$\sum_{\nu<l}\max(0,\theta-\nu/r_p)$ into a numerical lemma with explicit
positivity hypotheses. Split the argument into $\theta r_p\leq10$,
$10<\theta r_p<l$, and $l\leq\theta r_p$: the first uses the small-index
shortcut, while the others give the quadratic and linear lower bounds.
Use a natural-number floor and a finite range for the truncated sum,
and clear positive denominators before applying arithmetic tactics.}

    Combining this inequality with \eqref{roth:eq17} gives
    $$
    \min (\frac{1}{2}l \theta, \frac{1}{3}r_p \theta^2) \leq l (\Phi + \delta).
    $$
    Hence either $\frac{1}{2}l \theta \leq l (\Phi + \delta)$ or $\frac{1}{3}r_p \theta^2 \leq l (\Phi + \delta)$. In the first case, $\theta \leq 2(\Phi + \delta)$ which implies the inequality asserted in the lemma. In the second case,
        $$\frac{1}{3}r_p \theta^2 \leq l (\Phi + \delta) \leq (r_p + 1) (\Phi + \delta) < \frac{4}{3}r_p (\Phi + \delta),
        $$
        as $r_p > 10 \delta^{-1}$. These inequalities imply
        $$
        \theta < 2(\Phi + \delta)^{1/2} \leq 2(\Phi^{1/2} + \delta^{1/2}).
        $$
        This completes the proof of the lemma.
\end{proof}

\begin{lem}\label{roth:lem7}
    Let $m$ be a positive integer and let $\delta$ satisfy
    \begin{equation}
        \label{roth:eq19}
        0 < \delta < m^{-1}.
    \end{equation}
    Let $r_1, \dots, r_m$ be positive integers satisfying
    \begin{equation}
        \label{roth:eq20}
r_m > 10\delta^{-1}, \quad r_{j-1}/r_j > \delta^{-1}, \quad 2 \leq j \leq m.
        \end{equation}
        Let $q_1, \dots, q_m$ be positive integers satisfying
        \begin{equation}
        \label{roth:eq21}
\log q_1 > \delta^{-1} m (2m+1),
        \end{equation}
        \begin{equation}
        \label{roth:eq22}
r_j \log q_j \geq r_1 \log q_1, \quad \text{ for } 2 \leq j \leq m.
        \end{equation}
        Then
        \begin{equation}
        \label{roth:eq23}
\Theta_m(q_1^{\delta r_1}; q_1, \dots, q_m; r_1, \dots, r_m) < 10^m \delta^{(1/2)^m}.
        \end{equation}
\end{lem}
\begin{proof}
    Conditions \eqref{roth:eq21}--\eqref{roth:eq22} imply $q_j>1$ for every $j$. We will prove this lemma by induction on $m$.

    If $m=1$, Lemma \ref{roth:lem5} says
    $$
    \Theta_1(q_1^{\delta r_1}; q_1; r_1) \leq \frac{\delta r_1 \log q_1}{r_1 \log q_1}=\delta<10\delta^{1/2}
    $$

    Next, suppose $p \geq 2$ is an integer and that the lemma is true for $m= p-1$. We apply Lemma \ref{roth:lem6}. In order to do this we will need to estimate $M$ and $\Phi$ of Lemma \ref{roth:lem6}.

    We have
    $$
    M = (r_1 + 1)^{pl} l! 2^{plr_1} q_1^{\delta l r_1} \leq \left((r_1 + 1)^{p} l 2^{pr_1} q_1^{\delta r_1}\right)^l.
    $$
    Next, $l \leq r_p+1 < r_1 + 1 \leq 2^{r_1}$. Consequently,
    $$
    M < \left(2^{(2p+1)r_1}q_1^{\delta r_1} \right)^l < \left(e^{(2p+1)r_1}q_1^{\delta r_1} \right)^l.
    $$
    From Equation \ref{roth:eq21} for $m = p$, we have $2p+1 < \delta p^{-1} \log q_1$, and this gives
    $$
    M < \left(e^{(\delta p^{-1} \log q_1)r_1}q_1^{\delta r_1} \right)^l= \left(q_1^{(\delta p^{-1} )r_1}q_1^{\delta r_1} \right)^l= q_1^{\delta_1 l r_1}
    $$
    with $\delta_1 = \delta(1+p^{-1})$. Hence,
    $$
    \Theta_1(M; q_p; lr_p) \leq \Theta_1(q_1^{\delta_1 l r_1}; q_p; l r_p) \leq \frac{\log (q_1^{\delta_1 l r_1})}{l r_p \log q_p} \leq \frac{\delta_1 l r_1 \log q_1}{l r_1 \log q_1} = \delta_1,
    $$
    by Lemma \ref{roth:lem5}. Next,
    $$
    \Theta_{p-1}(M; q_1, \dots, q_{p-1}; lr_1, \dots, lr_{p-1}) \leq \Theta_{p-1}(q_1^{\delta_1 l r_1}; q_1, \dots, q_{p-1}; lr_1, \dots, lr_{p-1}).
    $$
    We verify the induction hypotheses for dimension $p-1$, parameter
    $\delta_1=\delta(1+1/p)$, and degrees $lr_1,\dots,lr_{p-1}$.
    First,
    \[
    0<(p-1)\delta_1<\frac{(p-1)(p+1)}{p^2}<1.
    \]
    Since $\delta_1>\delta$ and $l\geq1$, the degree-ratio and final-degree
    bounds follow from \eqref{roth:eq20}. Moreover,
    \[
    \log q_1>\delta^{-1}p(2p+1)
       >\delta_1^{-1}(p-1)(2p-1).
    \]
    The balanced-logarithm inequalities are unchanged when every degree is
    multiplied by $l$. Hence all four hypotheses hold. By the induction hypothesis,
    $$
    \Theta_{p-1}(q_1^{\delta_1 l r_1}; q_1, \dots, q_{p-1}; lr_1, \dots, lr_{p-1}) < 10^{p-1} \delta_1^{(1/2)^{p-1}}.
    $$
    Then
    $$
    \Phi < \delta_1 +  10^{p-1} \delta_1^{(1/2)^{p-1}} = \delta (1+p^{-1}) + 10^{p-1} \delta^{(1/2)^{p-1}} (1+p^{-1})^{(1/2)^{p-1}}
    $$
    $$
    < 2 \delta + 2 \cdot 10^{p-1} \delta^{(1/2)^{p-1}} < 3\cdot 10^{p-1} \delta^{(1/2)^{p-1}}.
    $$
    Here we have used the fact that $\delta_1 < 2 \delta$ and
    $$
    (1+p^{-1})^{(1/2)^{p-1}} < 2,
    $$
    which are both easily verified. Also $0<\delta<1$, so
    $\delta\leq\delta^{2^{-(p-1)}}$, justifying the last absorption above.
    For every allowed $l$, Lemma \ref{roth:lem6} therefore gives the same
    strict upper bound, and the maximum is over a nonempty finite set.
    Set $a=\delta^{2^{-p}}$. Then
    $\delta^{2^{-(p-1)}}=a^2\leq a$ and $\delta^{1/2}\leq a$. Consequently,
    \begin{align*}
    \Theta_p(q_1^{\delta r_1};q_1,\dots,q_p;r_1,\dots,r_p)
    &<2\left(3\cdot10^{p-1}a^2+
       \sqrt3\,10^{(p-1)/2}a+\delta^{1/2}\right)\\
    &\leq2\left(\frac3{10}+\frac{\sqrt3}{10^{(p+1)/2}}
                         +\frac1{10^p}\right)10^pa\\
    &\leq2\left(\frac3{10}+\frac{\sqrt3}{10^{3/2}}
                         +\frac1{10^2}\right)10^pa
      <10^pa.
    \end{align*}
    Here $p\geq2$ and the numerical coefficient in the last line is less
    than $1$ (for example, $\sqrt3/10^{3/2}<1/10$).
    This establishes the induction step, and we are done.
\end{proof}

\lean{For fixed $0<\delta<1$, define $a_0=\delta$ and
$a_{j+1}=\sqrt{a_j}$, and prove once that $a_j=\delta^{2^{-j}}$.
This lets the numerical induction use positivity and
$a_{j+1}^2=a_j$ instead of repeated real-exponent calculations.
Keep the maximum over $l$ explicitly finite: strict bounds on each
member imply a strict bound on a finite maximum, but not on an
arbitrary supremum.}

\section{Lattice count estimate}

\begin{lem}\label{roth:lem8}
    If $m\geq1$, $r_1, \dots, r_m$ are positive integers and $\lambda>0$, then the number of $m$-tuples of integers $(j_1, \dots, j_m)$ that satisfy
    $$
    0 \leq j_1 \leq r_1, \dots, 0 \leq j_m \leq r_m, \quad \frac{j_1}{r_1}+ \cdots + \frac{j_m}{r_m} \leq
\frac{m-\lambda}{2}
    $$
    does not exceed
    $$
    2 m^{1/2} \lambda^{-1} (r_1 + 1) \cdots (r_m+1).
    $$
\end{lem}
\begin{proof}
We will prove the result by induction on $m$. If $m=1$, we need to estimate the number of integers $j_1$ satisfying
$$
0 \leq j_1 \leq r_1, \quad j_1 \leq \frac{1}{2}(1-\lambda)r_1.
$$
If $\lambda > 1$ then we have no such $j_1$. If $0<\lambda\leq1$, then the number of $j_1$ from the first inequality is at most $r_1 + 1$, and this number is less than $2 \lambda^{-1}(r_1+1)$, establishing the case of $m=1$.

Now suppose $m >1$ and assume that the result is true for $m-1$. If $\lambda \leq 2m^{1/2}$ then the result is trivial, for in this case
$$
2m^{1/2}\lambda^{-1} (r_1 +1)\cdots (r_m+1) \geq (r_1 +1)\cdots (r_m+1)
$$
and the latter is equal to the number of solutions of $(j_1, \dots, j_m)$ satisfying
$$
    0 \leq j_1 \leq r_1, \dots, 0 \leq j_m \leq r_m.
    $$

So we assume $\lambda > 2 m^{1/2}$. For a fixed value of $j_m$, the induction hypothesis is applied to $j_1, \dots, j_{m-1}$ in the following way:
$$
    0 \leq j_1 \leq r_1, \dots, 0 \leq j_{m-1} \leq r_{m-1}, \quad \frac{j_1}{r_1}+ \cdots + \frac{j_{m-1}}{r_{m-1}} \leq
\frac{m-1-\lambda'}{2}
    $$
    where $\lambda'$ is chosen so that
    $$
    \frac{m-1-\lambda'}{2}=\frac{m-\lambda}{2}-\frac{j_m}{r_m},
    $$
    i.e., $\lambda' = \lambda - 1 + 2 j_m/r_m$. We note that $\lambda' > 0$. This is because
    $$
    \lambda' = \lambda - 1 + 2 j_m/r_m > 2m^{1/2}-1 > 0.
    $$
    By the induction hypothesis, the number of $m-1$-tuples is at most
    $$
    2(m-1)^{1/2} (\lambda')^{-1} (r_1+1)\cdots(r_{m-1}+1).
    $$
    Summing these bounds over $0\leq j_m\leq r_m$, it is enough to show
\[
2\sqrt{m-1}\prod_{i=1}^{m-1}(r_i+1)
 \sum_{j_m=0}^{r_m}\frac1{\lambda-1+2j_m/r_m}
 \leq\frac{2\sqrt m}{\lambda}\prod_{i=1}^{m}(r_i+1).
\]
To finish the proof of the lemma we prove the following statement: Let $r\geq1$ and $m\geq2$ be integers and suppose $\lambda>2\sqrt m$. Then
$$
\sum_{j=0}^r (\lambda -1 + 2j/r)^{-1} < \lambda^{-1} (m-1)^{-1/2} m^{1/2} (r+1).
$$
There are two cases. If $r$ is even, we replace $j$ by $\frac{1}{2}r+k$ to obtain the sum
$$
\sum_{k=-r/2}^{r/2} (\lambda + 2k/r)^{-1} = \lambda^{-1} + \sum_{k=1}^{r/2} 2\lambda (\lambda^2 - 4k^2/r^2)^{-1}
$$
$$
\leq \lambda^{-1} + \sum_{k=1}^{r/2} 2\lambda (\lambda^2-1)^{-1}
$$
$$
\leq (r+1)\lambda^{-1}(1-\lambda^{-2})^{-1}.
$$
Now
$1-\lambda^{-2}>1-1/(4m)>\sqrt{1-1/m}>0$.
For the middle inequality, square the positive quantities and use
$(1-1/(4m))^2-(1-1/m)=1/(2m)+1/(16m^2)>0$.
Taking reciprocals gives the result for even $r$.
If $r$ is odd, we obtain
$$
\sum_{k=0}^{(r-1)/2} \left(\frac{1}{\lambda - \frac{1+2k}{r}}+  \frac{1}{\lambda + \frac{1+2k}{r}}\right)
= \sum_{k=0}^{(r-1)/2}
\frac{2\lambda}{\lambda^2 - \left(\frac{1+2k}{r}\right)^2}
$$
$$
\leq \sum_{k=0}^{(r-1)/2}
\frac{2\lambda}{\lambda^2 - 1}= \frac{r+1}{2}\cdot \frac{2\lambda}{\lambda^2 - 1} = \frac{\lambda^{-1}(r+1)}{1-\lambda^{-2}}.
$$
The rest of the argument is as before.
\end{proof}

\lean{Represent the lattice box by functions whose $j$th value lies in
\texttt{Fin}$(r_j+1)$, then filter by the displayed weight inequality.
Its cardinality is the product $\prod_j(r_j+1)$, and induction counts the
fibres of the last-coordinate projection. Prove the reciprocal-sum
estimate separately by pairing $j$ with $r-j$, keeping the finite counting
argument distinct from the positivity and square-root inequalities.}

\section{The main lemma: Roth's Lemma 9}
Let $\alpha$ be an irrational real algebraic number and let $\kappa\in\R$. Suppose there are infinitely many pairs of integers $h,q$ with $q>0$ such that
\begin{equation}\label{roth:eq1}
\left| \alpha - \frac{h}{q}\right| <\frac{1}{q^\kappa}.
\end{equation}
\begin{thm}[Roth]
    Under the above conditions, $\kappa \leq 2$.
\end{thm}
\begin{lem}
    If the theorem fails, then it also fails for an irrational real algebraic integer.
\end{lem}
\begin{proof}
Suppose the theorem fails for a real irrational algebraic number $\alpha$: there is a $\kappa>2$ such that Equation \eqref{roth:eq1} holds for infinitely many integer pairs $h,q$ with $q>0$.

\begin{lem}
    If $\alpha$ is an algebraic number, there is an integer $M>0$ such that $M\alpha$ is an algebraic integer.
\end{lem}
\begin{proof}
    Clear denominators in a nonzero rational polynomial annihilating
    $\alpha$ and change its sign so that its leading coefficient is a
    positive integer $M$. Write it as
    $MX^d+b_{d-1}X^{d-1}+\cdots+b_0$, where $d\geq1$.
    Then $\beta=M\alpha$ satisfies the monic integral equation
    \[
    \beta^d+\sum_{j=0}^{d-1}b_jM^{d-1-j}\beta^j=0.
    \]
    Thus $M\alpha$ is an algebraic integer.
\end{proof}
Let $\alpha' = M \alpha$. Then multiplying Equation \eqref{roth:eq1} by $M$ gives
$$
\left| \alpha' - \frac{h'}{q}\right| <\frac{M}{q^\kappa}
$$
with $h' = Mh$. Pick any $\kappa'$ such that $2 < \kappa' < \kappa$. Then for $q$ large enough,
$$
\frac{M}{q^\kappa} < \frac{1}{q^{\kappa'}}.
$$
For a fixed upper bound on $q$, there are only finitely many possible
integers $h$ satisfying the approximation inequality. The infinitely many
original solutions therefore have unbounded denominators, and discarding
bounded $q$ leaves infinitely many. For such $q$ we have
$$
\left| \alpha' - \frac{h'}{q}\right| <\frac{1}{q^{\kappa'}}.
$$
The algebraic integer $\alpha'$ is still irrational, because $M$ is a
nonzero integer. It has infinitely many approximations with exponent
$\kappa'>2$, as required.
\end{proof}

\lean{Separate the algebraic lemma producing a positive integer $M$
with $M\alpha$ integral over $\Z$ from the estimate
$M q^{-\kappa}<q^{-\kappa'}$ for all sufficiently large $q$, where
$2<\kappa'<\kappa$. Also prove separately that solutions with bounded
denominator form a finite set. Removing that finite set then preserves
infinitely many approximations, without mixing the integrality proof
with the real-power and infinitude arguments.}

Let $\alpha$ be an irrational real algebraic integer such that
$$
\left| \alpha - \frac{h}{q}\right| <\frac{1}{q^{\kappa}}
$$
for infinitely many pairs of integers $h,q$ with $q>0$ and some $\kappa>0$. Our goal is to prove that $\kappa \leq 2$.

Since $\alpha$ is an algebraic integer, by definition, there is a polynomial
$$
f(x) = x^n + a_1 x^{n-1} + \cdots + a_n
$$
with integral coefficients such that $f(\alpha)=0$. Its degree $n$ is at least $2$, since $\alpha$ is irrational. For future reference we let
\begin{equation}
    \label{roth:eq28}
    A = \max(1, |a_1|, \dots, |a_n|).
\end{equation}

\begin{lem}\label{lem:m-delta}
    For every integer $m>16n^2$ there is a real $\delta>0$ such that
    \begin{equation}
        \label{roth:eq29}
0 < \delta < m^{-1}
    \end{equation}
    and
    \begin{equation}
        \label{roth:eq30}
        10^m \delta^{(1/2)^m} + 2(1+3 \delta) n m^{1/2} < \frac{1}{2}m.
    \end{equation}
\end{lem}
\begin{proof}
Write the second inequality as
\[
2n\sqrt m+
\left(10^m\delta^{2^{-m}}+6\delta n\sqrt m\right)
<\frac m2.
\]
For fixed $m$, the expression in parentheses tends to $0$ as $\delta\to 0^+$. Thus the desired $\delta$ can exist precisely when
\[
2n\sqrt m<\frac m2,
\]
or equivalently,
\[
m>16n^2.
\]

More generally, for any fixed integer $m>16n^2$, put
$E=m/2-2n\sqrt m>0$ and choose
\[
\delta=\frac12\min\left\{\frac1m,\frac{E}{12n\sqrt m},
              \left(\frac{E}{2\cdot10^m}\right)^{2^m}\right\}.
\]
All entries are positive. This gives
$10^m\delta^{2^{-m}}<E/2$ and $6\delta n\sqrt m<E/2$, proving the
claim for the fixed $m$. Every smaller positive $\delta$ also works,
since the left-hand side of \eqref{roth:eq30} is strictly increasing in
$\delta$. A convenient explicit choice of $m$ is
\[
\boxed{m=(4n+1)^2.}
\]
Then $\sqrt m=4n+1$, and hence
\[
\frac m2-2n\sqrt m
=\frac{\sqrt m}{2}.
\]

Now choose
\[
\boxed{
\delta=
\frac12\min\left\{
\frac1m,
\frac1{24n},
\left(\frac{\sqrt m}{4\cdot 10^m}\right)^{2^m}
\right\}.
}
\]

Indeed, the first entry in the minimum gives
\[
0<\delta<\frac 1 m.
\]
The third entry gives
\[
\delta^{(1/2)^{m}} < \frac{\sqrt m}{4\cdot10^m},
\]
and therefore
\[
10^m\delta^{2^{-m}}<\frac{\sqrt m}{4}.
\]
The second entry gives
\[
6\delta n\sqrt m<\frac{\sqrt m}{4}.
\]
Consequently,

\begin{align*}
10^m \delta^{(1/2)^{m}} +2(1+3\delta)n\sqrt m
&=
2n\sqrt m+
10^m\delta^{(1/2)^{m}}
+6\delta n\sqrt m \\
&<
2n\sqrt m+\frac{\sqrt m}{2} \\
&=\frac m2.
\end{align*}

Thus this explicit pair $(m,\delta)$ satisfies both required inequalities,
and the preceding construction proves the stronger assertion for each
fixed $m>16n^2$.
\end{proof}

\lean{For each fixed $m>16n^2$, prove the stronger statement that some
$\delta_0>0$ makes \eqref{roth:eq29} and \eqref{roth:eq30} hold for every
$0<\delta<\delta_0$. The explicit positive-minimum construction above
can provide this threshold. The final theorem must choose $m$ using
$\kappa$ and may then shrink $\delta$ further, so do not fix $m$ globally
to the special example $(4n+1)^2$.}

Let $m$ and $\delta$ be as in the lemma. We choose positive integers $r_1, \dots, r_m$ such that
\begin{equation}
    \label{roth:eq31}
    r_m > 10 \delta^{-1}, \quad r_{j-1}/r_j > \delta^{-1}, \quad 2 \leq j \leq m.
\end{equation}
This is easily arranged. Next, we pick a positive integer $q_1>1$ such that
\begin{equation}
    \label{roth:eq32}
    \delta^2 \log q_1 > 2m+1 + 2m \log (1+A) + 2m \log (1 + |\alpha|),
\end{equation}
with $A$ as in Equation \eqref{roth:eq28}.

Next, we pick positive integers $q_2, \dots, q_m$ such that
\begin{equation}
    \label{roth:eq33}
    r_j \log q_j \geq r_1 \log q_1, \quad 2 \leq j \leq m.
\end{equation}

\begin{lem}\label{lem:lem7-satisfied}
    If $m, \delta, r_1, \dots, r_m, q_1, \dots, q_m$ satisfy Equations \eqref{roth:eq29}-\eqref{roth:eq33} then they satisfy the conditions of Lemma \ref{roth:lem7}.
\end{lem}
\begin{proof}
    We just need to verify $\delta \log q_1 > m (2m+1)$ but this is immediate from \eqref{roth:eq29} and \eqref{roth:eq32}.
\end{proof}

Define $\lambda, \gamma, \eta, B_1$ by
\begin{equation}
    \label{roth:eq34}
    \lambda = 4(1+3 \delta)nm^{1/2},
\end{equation}
\begin{equation}
    \label{roth:eq35}
    \gamma = \frac{1}{2}(m - \lambda),
\end{equation}
\begin{equation}
    \label{roth:eq36}
    \eta = 10^m \delta^{(1/2)^m},
\end{equation}
\begin{equation}
    \label{roth:eq37}
    B_1= [q_1^{\delta r_1}].
\end{equation}

Inequality \eqref{roth:eq30} is equivalent to
\begin{equation}
    \label{roth:eq38}
    \eta < \gamma.
\end{equation}
We also note that $B_1$ is a \emph{large} integer, because
$$
B_1 > q_1^{\delta r_1} -1 = (e^{\delta \log q_1})^{r_1} -1 > (e^{m(2m+1)})^{r_1}-1 >(e^{16n^2(2 \times 16 n^2+1)})^{r_1}-1
$$
$$
\geq (e^{16 \times 2^2(32 \times 2^2+1)})^{r_1}-1 = e^{8256r_1} -1 > e^{82560}-1.
$$
In particular, writing $X=q_1^{\delta r_1}$, we have $X>4$, whence
$B_1=\lfloor X\rfloor>X-1>\sqrt X=q_1^{\delta r_1/2}$.
Thus
\[
\log B_1>\tfrac12\delta r_1\log q_1.
\]
Together with \eqref{roth:eq32}, this gives the three bounds used below:
\begin{equation}\label{roth:three-small-factors}
2^{mr_1}<B_1^\delta,\qquad
(1+A)^{mr_1}<B_1^\delta,\qquad
(1+|\alpha|)^{mr_1}<B_1^\delta.
\end{equation}
For the first, use
$m\log2<\tfrac12\delta^2\log q_1$; for the other two, use the
corresponding logarithmic terms in \eqref{roth:eq32}.
\lean{Put $X=q_1^{\delta r_1}$ and isolate the floor estimate
$\sqrt X<\lfloor X\rfloor\leq X$ for $X>4$.
Use it to obtain $\log B_1>\tfrac12\delta r_1\log q_1$, then prove once that
$2^{mr_1}$, $(1+A)^{mr_1}$ and $(1+|\alpha|)^{mr_1}$ are all less than
$B_1^\delta$, with all positive-base hypotheses explicit.
Retain the non-strict upper bound $B_1\leq X$: equality is possible when
$X$ is an integer.}

Now we come to Roth's main lemma.

\begin{lem}
\label{roth:lem9}
    Suppose \eqref{roth:eq29}-\eqref{roth:eq33} are satisfied, and suppose $h_1, \dots, h_m$ are integers that are coprime to $q_1, \dots, q_m$ respectively. Then there exists a nonzero polynomial $Q(x_1, \dots, x_m)$ with integral coefficients, of degree at most $r_j$ in $x_j$ for $1 \leq j \leq m$, such that
\begin{enumerate}
    \item The index of $Q$ at the point $(\alpha, \dots, \alpha)$ relative to $r_1, \dots, r_m$ is at least $\gamma - \eta >0$;
    \item $Q(h_1/q_1, \dots, h_m/q_m) \ne 0$;
    \item for all non-negative integers $i_1, \dots, i_m$ we have
    $$
    \left| \Delta_{i_1, \dots, i_m} Q(\alpha, \dots, \alpha) \right| < B_1^{1+3 \delta}.
    $$
\end{enumerate}
\end{lem}
\begin{proof}
    Consider the set $\mathcal W$ of all polynomials
    $$
    W(x_1, \dots, x_m) =\sum_{s_1=0}^{r_1} \dots \sum_{s_m=0}^{r_m} c(s_1, \dots, s_m) x_1^{s_1} \cdots x_m^{s_m}
    $$
    where the coefficients $c(s_1, \dots, s_m)$ are integers satisfying
    $$
    0 \leq c(s_1, \dots, s_m) \leq B_1.
    $$
    The number of such polynomials is
    $$
    |\mathcal W| =  (B_1+1)^r
    $$
    with
     \begin{equation}
     \label{roth:eq43}
     r = (r_1 + 1) \cdots (r_m+1).
 \end{equation}

 For $W \in \mathcal W$ consider the derivatives
 $$
 \Delta_{j_1, \dots, j_m} W
 $$
 for integers $j_1, \dots, j_m$ satisfying
  \begin{equation}
     \label{roth:eq44}
 0 \leq j_1 \leq r_1, \dots, 0 \leq j_m \leq r_m, \quad \frac{j_1}{r_1}+ \cdots + \frac{j_m}{r_m} \leq \gamma.
 \end{equation}
 By Lemma \ref{roth:lem8} and Equation \eqref{roth:eq35}, the number $D$ of such derivatives satisfies
 \begin{equation}
     \label{roth:eq45}
     D \leq 2m^{1/2} \lambda^{-1} r,
 \end{equation}
 with $r$ as in \eqref{roth:eq43}. For each such derivative let
 $$
 T_{j_1, \dots, j_m}(W; x)
 $$
 be the remainder of the division of $\Delta_{j_1, \dots, j_m}W(x, \dots, x)$ by $f(x)$.  Since $f$ is a monic polynomial, $T_{j_1, \dots, j_m}(W; x)$ will have integer coefficients, and its degree is less than or equal to $n-1$.

Our first task is to find an upper bound for the sizes of the coefficients of the polynomials $T_{j_1, \dots, j_m}(W; x)$. For $W \in \mathcal W$ we have
$$
abc(\Delta_{j_1, \dots, j_m} W) \leq 2^{r_1 +\dots + r_m} B_1 \leq 2^{mr_1} B_1 < B_1^{1+\delta}
$$
by \eqref{roth:three-small-factors}. Indeed, \eqref{roth:eq32} gives
$mr_1\log2<\tfrac12\delta^2r_1\log q_1<\delta\log B_1$.
When we replace $x_1, \dots, x_m$ by the same variable $x$ the absolute value of the coefficients of the resulting polynomial is at most the number of terms of the polynomial, which is at most $r$, multiplied by the coefficient bound $B_1^{1+\delta}$. Now,
$$
r = (r_1 + 1) \cdots (r_m + 1) \leq 2^{r_1 + \dots + r_m} \leq 2^{mr_1} < B_1^\delta,
$$
and this implies that $rB_1^{1+ \delta} < B_1^{1 + 2 \delta}$.  Consequently, the coefficients of $\Delta_{j_1, \dots, j_m}W (x, \dots, x)$ are bounded in absolute value by $B_1^{1+ 2\delta}$. Next we examine the effect of division by $f$ on these coefficients. Write
$$
\Delta_{j_1, \dots, j_m}W (x, \dots, x) = w_s x^s + w_{s-1} x^{s-1} + \cdots + w_0.
$$
When we divide this polynomial by $f$, the first operation we do (assuming $s \geq n$) is to subtract $w_s x^{s-n} f(x)$ and this gives a new polynomial whose coefficients are either of the form $w_\nu-a_{s-\nu}w_s$ (when $1\leq s-\nu\leq n$) or of the form $w_\nu$.  Since the $a_{s-\nu}$'s
are bounded by $A$, the absolute values of the resulting coefficients are at most $$(1 + A)  \cdot abc\, (\Delta_{j_1, \dots, j_m}W (x, \dots, x)) \leq (1+A) B_1^{1 + 2 \delta}.$$
If the input polynomial is zero, its remainder is zero. Otherwise the
number of leading-term eliminations is at most $\max(0,s-n+1)$: there are
none when $s<n$, and each elimination lowers the degree by at least one.
It follows that
\[
abc(T_{j_1,\dots,j_m}(W;x))
 \leq(1+A)^{\max(0,s-n+1)}B_1^{1+2\delta}
 \leq(1+A)^{mr_1}B_1^{1+2\delta}<B_1^{1+3\delta},
\]
using $s\leq\sum_jr_j\leq mr_1$ and
\eqref{roth:three-small-factors}.
\lean{Use
\texttt{Polynomial.modByMonic}\footnote{\url{https://leanprover-community.github.io/mathlib4_docs/Mathlib/Algebra/Polynomial/Div.html}}
directly in $\Z[X]$, rather than dividing over $\Q$ and later recovering
integrality. Prove the coefficient bound by induction on degree, treating
zero and degree less than $n$ first; each leading-term elimination
multiplies the bound by at most $1+A$ and lowers the degree.
Handle diagonal substitution separately, since identifying variables
can merge coefficients.}

Pad each remainder with zero coefficients to length $n$. Each integer coefficient has fewer than $1+2B_1^{1+3\delta}$ possible values. Since there are $D$ derivatives, the number of possible remainder tuples is at most
$$
(1+ 2B_1^{1+3 \delta})^{nD}.
$$
By \eqref{roth:eq45} and the definition of $\lambda$ in \eqref{roth:eq34} we have
$$
(1 + 3 \delta) n D \leq 2(1+3 \delta) nm^{1/2} \lambda^{-1} r = \frac{1}{2}r.
$$
Put $c=1+3\delta>1$. Since $B_1>1$,
\[
1+2B_1^c<[2(1+B_1)]^c.
\]
Indeed, $(1+B_1)^c\geq1+B_1^c$ and $2^c\geq2$, so the right-hand side
is at least $2+2B_1^c$. The tuple $j=0$ is allowed because $\gamma>0$,
so $D\geq1$. Using $cnD\leq r/2$, we obtain
\begin{align*}
(1+2B_1^c)^{nD}
 &<[2(1+B_1)]^{cnD}\leq[2(1+B_1)]^{r/2}\\
 &<(1+B_1)^r=|\mathcal W|.
\end{align*}
The last strict inequality follows from $2(1+B_1)<(1+B_1)^2$, or
$B_1>1$. Thus there are fewer remainder tuples than elements of $\mathcal W$.
 By the pigeonhole principle, there are two distinct polynomials $W',W''\in\mathcal W$ such that
$$
\Delta_{j_1, \dots, j_m} W' (x, \dots, x) - \Delta_{j_1, \dots, j_m} W'' (x, \dots, x)
$$
is divisible by $f(x)$ for all $j_1, \dots, j_m$ that satisfy \eqref{roth:eq44}. Let $W^*=W'-W''$. Hasse differentiation, diagonal substitution and
remainder modulo the fixed monic $f$ are additive, so equality of remainder
tuples implies the stated divisibility for the difference. For all $j_1, \dots, j_m$ satisfying \eqref{roth:eq44}
$$
\Delta_{j_1, \dots, j_m} W^* (x_1, \dots, x_m)
$$
are zero when $x_1 = \dots = x_m = \alpha$. Multi-indices outside the degree box give identically zero derivatives as well. Hence the index of $W^*$ at the point $(\alpha, \dots, \alpha)$ relative to $r_1, \dots, r_m$ is at least $\gamma$. Also $W^*$ is not identically zero, and its coefficients in absolute value do not exceed $B_1$.

\lean{Represent each $W$ by a coefficient array on the finite monomial
box, with entries in $\{0,\ldots,B_1\}$, and map it to the $nD$ integer
coefficients of its remainder tuple. Apply
finite pigeonhole\footnote{\url{https://leanprover-community.github.io/mathlib4_docs/Mathlib/Combinatorics/Pigeonhole.html}}
using the strict cardinality estimate above. Distinct arrays give a
nonzero difference with coefficient bound $B_1$, not $2B_1$, because both
arrays lie in $[0,B_1]$ coordinatewise; equality of remainders then forces
the required low-weight derivatives to vanish at $\alpha$.}

This means that
$$
W^* (x_1, \dots, x_m) \in \mathcal R_m(q_1^{\delta r_1}; r_1, \dots, r_m).
$$
We noted in Lemma \ref{lem:lem7-satisfied} that the conditions of Lemma \ref{roth:lem7} are satisfied. By Lemma \ref{roth:lem7} the index of $W^*$ at $(h_1/q_1, \dots, h_m/q_m)$ relative to $r_1, \dots, r_m$ is less than $\eta$ defined in \eqref{roth:eq36}. Hence $W^*$ has some derivative
$$
Q(x_1, \dots, x_m) = \Delta_{k_1, \dots, k_m} W^*
$$
with $$
\frac{k_1}{r_1}+ \cdots + \frac{k_m}{r_m} < \eta
$$
such that
$$
Q(h_1/q_1, \dots, h_m/q_m) \ne 0.
$$
By Lemma \ref{lem:index-of-derivative} the index of $Q$ at the point $(\alpha, \dots, \alpha)$ relative to $r_1, \dots, r_m$ is at least $\gamma - \eta$. This $Q$ has properties (1) and (2) of Lemma \ref{roth:lem9}. Since the coefficients of $W^*$ have absolute value at most $B_1$, the coefficients of $Q$ have absolute value at most
$$
2^{r_1 + \dots + r_m} B_1 \leq 2^{mr_1} B_1 < B_1^{1 + \delta}.
$$
Hence the coefficients of any further derivative of $Q$ have absolute value at most
$$
2^{r_1 + \dots + r_m}B_1^{1+ \delta} <B_1^{1+ 2 \delta}.
$$
To justify evaluation, for every $t\geq0$ and integer $r\geq0$,
$\sum_{j=0}^rt^j\leq(1+t)^r$, since each binomial coefficient
$\binom rj$ is at least $1$. Thus the coefficient bound gives
\[
|(\Delta_iQ)(\alpha,\dots,\alpha)|
 \leq abc(\Delta_iQ)\prod_{j=1}^m\sum_{s=0}^{r_j}|\alpha|^s.
\]
For a zero derivative the desired estimate is immediate; otherwise the
strict coefficient bound yields
$$
|\Delta_{i_1, \dots, i_m} Q(\alpha, \dots, \alpha)| < B_1^{1+ 2 \delta} (1 + |\alpha|)^{r_1 + \dots + r_m},
$$
and this implies (3), since \eqref{roth:three-small-factors} gives
$$
(1 + |\alpha|)^{mr_1} < B_1^\delta.
$$
This completes the proof of Lemma \ref{roth:lem9}.
\lean{Keep differentiation, diagonal substitution and point evaluation
as three separate estimates. A fixed Hasse derivative has no coefficient
collisions; diagonal substitution may merge monomials and needs the
factor $\prod_j(r_j+1)$. For evaluation, sum the absolute coefficients
against powers of $|\alpha|$ and use
$\sum_{s=0}^r t^s\leq(1+t)^r$ for $t\geq0$; this absorbs the monomial
sum without an additional box-cardinality factor.}

\end{proof}
\section{The proof of Roth's theorem}
In this section we complete the proof of Roth's theorem using Lemma \ref{roth:lem9}. Let $\kappa > 2$ and suppose that there are infinitely many pairs of integers $h, q$, $q>0$, such that
\begin{equation}
    \label{roth:eq46}
    \left|\alpha - \frac{h}{q} \right|< \frac{1}{q^\kappa}.
\end{equation}
We will derive a contradiction.
We first justify choosing reduced solutions of arbitrarily large denominator.
Reducing $h/q$ to $a/b$ with $b>0$ and $\gcd(a,b)=1$ preserves the
inequality, since $b\leq q$ and $\kappa>0$.
For a fixed rational value $a/b$, irrationality gives
$\varepsilon=|\alpha-a/b|>0$. The inequality
$\varepsilon<q^{-\kappa}$ bounds its original denominators $q$, so only
finitely many original pairs can represent that rational value.
There are therefore infinitely many distinct reduced solutions.
For bounded $b$, their numerators $a$ lie in finite intervals, so the
reduced denominators are unbounded as well.
We choose a positive integer $m$ large enough that $m > 4nm^{1/2}$
and
\begin{equation}
    \label{roth:eq47}
\frac{2m}{m - 4nm^{1/2}} < \kappa,
\end{equation}
which is possible as $\kappa > 2$. Let $\eta$ be as in \eqref{roth:eq36}. For $\delta$ small enough we have
$$
m - 4(1 + 3 \delta) nm^{1/2} - 2 \eta >0.
$$
This condition is the same as \eqref{roth:eq30}. By the fixed-$m$ form of Lemma \ref{lem:m-delta}, we can choose $\delta$ so it also satisfies \eqref{roth:eq29}. Decreasing $\delta$ further if necessary, we also arrange that
\begin{equation}
    \label{roth:eq48}
    \frac{2m(1+4 \delta)}{m-4(1+3 \delta)nm^{1/2} - 2 \eta} < \kappa.
\end{equation}
This is possible because the left-hand side tends to the strictly smaller quantity in \eqref{roth:eq47} as $\delta\to0^+$. By the definition of $\gamma$ and $\eta$ from \eqref{roth:eq35} and \eqref{roth:eq36}, \eqref{roth:eq48} is equivalent to
\begin{equation}
    \label{roth:eq49}
    \frac{m(1+ 4 \delta)}{\gamma - \eta} < \kappa.
\end{equation}
Next, pick a reduced solution $h_1,q_1$ of \eqref{roth:eq46} such that $q_1$ satisfies \eqref{roth:eq32}. This means $q_1$ is very large. We then choose solutions $h_2, q_2; \dots; h_m, q_m$ with $\gcd(h_j,q_j)=1$ for all $j$ to satisfy
\begin{equation}
    \label{roth:eq50}
    \frac{\log q_j}{\log q_{j-1}}> \frac{2}{\delta}, \quad 2 \leq j \leq m.
\end{equation}
We now take $r_1$ to be any integer satisfying
\begin{equation}
    \label{roth:eq51}
    r_1 > \frac{10 \log q_m}{\delta \log q_1}
\end{equation}
and define integers $r_2, \dots, r_m$ by
\begin{equation}
    \label{roth:eq52}
    \frac{r_1 \log q_1}{\log q_j} \leq r_j < 1 + \frac{r_1 \log q_1}{\log q_j}, \quad 2 \leq j \leq m.
\end{equation}
Then \eqref{roth:eq33} is satisfied. Also,
\begin{equation}
    \label{roth:eq53}
    \frac{r_j \log q_j}{r_1 \log q_1} < 1+ \frac{\log q_j}{r_1\log q_1}\leq 1 + \frac{\log q_m}{r_1 \log q_1} < 1+ \frac{1}{10}\delta.
\end{equation}
The conditions \eqref{roth:eq31} are also satisfied since
$$
r_m \geq \frac{r_1 \log q_1}{\log q_m }> 10 \delta^{-1},
$$
and
$$
\frac{r_{j-1}}{r_j}> \frac{\log q_j}{\log q_{j-1}}(1+ \frac{1}{10} \delta)^{-1} > \delta^{-1}
$$
by \eqref{roth:eq52}, \eqref{roth:eq53}, and \eqref{roth:eq50}.

\lean{Start from unbounded reduced denominators and choose the finite
sequence of approximants satisfying \eqref{roth:eq50}. After choosing
$r_1$, define $r_j=\lceil r_1\log q_1/\log q_j\rceil$ for $j\geq2$,
and collect positivity, \eqref{roth:eq31}, \eqref{roth:eq33} and
\eqref{roth:eq53} in one rounding lemma. Use the
Diophantine approximation module\footnote{\url{https://leanprover-community.github.io/mathlib4_docs/Mathlib/NumberTheory/DiophantineApproximation/Basic.html}}
as a reference for the rational-approximation formulation, not as a
substitute for proving these selection and rounding steps.}

By Lemma \ref{roth:lem9} there exists a polynomial $Q(x_1, \dots, x_m)$ satisfying the assertions of that lemma. By (2) of Lemma \ref{roth:lem9}, $Q(h_1/q_1, \dots, h_m/q_m) \ne 0$. Also, since $Q$ has degree at most $r_j$ in $x_j$ for $1 \leq j \leq m$, the number
$\bigl(\prod_jq_j^{r_j}\bigr)Q(h_1/q_1,\dots,h_m/q_m)$ is a nonzero
integer: each monomial denominator divides $\prod_jq_j^{r_j}$.
Its absolute value is at least $1$, and therefore
\begin{equation}
    \label{roth:eq54}
    |Q(h_1/q_1, \dots, h_m/q_m)| \geq q_1^{-r_1} \cdots q_m^{-r_m} > q_1^{-mr_1(1+\delta)}
\end{equation}
by \eqref{roth:eq53}.  On the other hand,
we have
$$
Q(h_1/q_1, \dots, h_m/q_m) = \sum_{i_1=0}^{r_1} \cdots \sum_{i_m =0}^{r_m} \Delta_{i_1, \dots, i_m} Q(\alpha, \dots, \alpha) (h_1/q_1-\alpha)^{i_1} \cdots (h_m/q_m-\alpha)^{i_m}.
$$
By (1) of Lemma \ref{roth:lem9} the terms with
$$
\frac{i_1}{r_1}+ \cdots + \frac{i_m}{r_m} < \gamma - \eta
$$
vanish. In every other term, the weight is positive, so at least one $i_j$ is positive and the product estimate is strict. We have
$$
\left|  (h_1/q_1-\alpha)^{i_1} \cdots (h_m/q_m-\alpha)^{i_m} \right| < \frac{1}{(q_1^{i_1} \cdots q_m^{i_m})^\kappa} \leq q_1^{-r_1(\gamma - \eta)\kappa},
$$
since $q_j \geq q_1^{r_1/r_j}$ by \eqref{roth:eq52}. By using (3) of Lemma \ref{roth:lem9} we have
$$
\left| Q(h_1/q_1, \dots, h_m/q_m)\right| < (r_1 + 1) \cdots (r_m + 1) B_1^{1+3 \delta} q_1^{-r_1(\gamma - \eta)\kappa}
$$
$$
< B_1^{1+ 4 \delta} q_1^{-r_1(\gamma-\eta)\kappa}
$$
$$
\leq q_1^{(1+4\delta)\delta r_1-r_1(\gamma-\eta)\kappa}.
$$
Here the last comparison uses $B_1\leq q_1^{\delta r_1}$; the
preceding strict inequalities preserve a strict upper bound for $|Q|$.
Comparing this with \eqref{roth:eq54}, and using $q_1>1$, we obtain
$$
- m r_1(1+\delta) < (1+4 \delta) \delta r_1 - r_1(\gamma - \eta) \kappa,
$$
or
$$
\kappa < \frac{m(1+\delta) + \delta(1+4 \delta)}{\gamma - \eta} < \frac{m(1+4 \delta)}{\gamma - \eta},
$$
where the second inequality uses
$1+4\delta<3m$, which follows from $m>16n^2\geq64$ and $\delta<1/m$.
This contradicts \eqref{roth:eq49} and finishes the proof of the theorem.
\lean{Prove denominator clearing and the Taylor upper bound as separate lemmas.
With $q_1>1$ and $r_1>0$, compare logarithms and cancel $r_1\log q_1$.
For strictness, use $\gamma-\eta>0$ to ensure each surviving Taylor term
has a positive exponent; treat zero coefficients separately.}

\end{document}
